\section{Inputs to the pointwise lower law}
\label{app:v51-input-map}
The shortest new argument uses the following four inputs from the
retained actual-return theory. The source partition is
\eqref{eq:v51-exact-positive-split}; the protected correction is
Theorem~\ref{thm:v46-protected-correction}; the complete decision
correction is Corollary~\ref{cor:v48-decision-all-band}; and the
physical local-variation estimate is
Theorem~\ref{thm:v49-physical-remainder-local}. The fixed-band main
term is the full-source formula in the proof of
Theorem~\ref{thm:v43-arithmetic-raw-criterion}. None of these inputs
is a pointwise bound on the unprotected physical remainder.

For the inherited physical multiplier the following table fixes the
hypothesis-to-source correspondence. External lemma numbering follows
arXiv:1902.06850v1, as specified in \cite{DPZ}.

\smallskip
\noindent\begin{tabular}{p{0.31\linewidth}p{0.61\linewidth}}
Requirement & Location and convention in this article\\[3pt]
Finite partition complexity & Lemma~\ref{lem:v49-physical-envelope}: a fixed candidate-center set and bounded trigonometric degree on each homogeneity strip.\\[3pt]
Stable transversality & Equation~\eqref{eq:v49-transverse-derivatives}: clearance levels have positive slope, stable graphs negative slope.\\[3pt]
Boundary neighborhoods & The same lemma: an interval of width $s$ has length $O(s)$ on a stable graph; coalescing boundaries introduce no inverse width.\\[3pt]
Grazing thresholds & Two horizontal angle cuts; the original homogeneity strips remain in force. No inverse change from $p$ is used at grazing.\\[3pt]
Multiplier exponents & Equation~\eqref{eq:v35-exponent-choice} and \cite[Lemma 3.3(b)]{DPZ}; the indicator is constant on each partition cell.\\[3pt]
Weak gain and physical product & Proposition~\ref{prop:v49-physical-multiplier}: stable-length exponent $1/8$, with a uniform strong bound, not a small strong bound.\\
\end{tabular}
\smallskip

The preceding incident disk does not thicken the clearance envelope:
its physical incoming line intersects that disk and therefore has
line distance at most $R$, while the envelope requires at least $R$.
Their common boundary is the tangency condition. Clearance of flight
$j$ remains marked at collision $j+1$, and the occupation sum is
$\sum_{a=0}^{m-1}\eta_R\circ T_R^a$. The source factor $1/c$ is
paid when dropping the section restrictions in the positive upper
comparison, before the new argument begins.

These are identified continuum inputs, not claims of independent
specialist verification. In particular the matching geometry, the
strong-completion multiplication, the moving-peak interpolation and
the all-depth finite-count remainder still carry the audit scope
described in the relevant proofs.

\paragraph{Comparison of theorem statements.}
The cell-index mixing local theorem with endpoint observables and
scatterer perturbations in \cite{DPZ}, and the abstract suspension
local theorem in \cite{DN}, remain the relevant comparison points.
The new conclusion here is neither a differentiation of their interval
laws nor an assertion that a singleton lattice event is a new topology.
It is the lower inequality for the actual four-coordinate return
density, obtained from the positive source partition above. The abstract
real-analysis proposition \ref{prop:v51-positive-principle} makes
this inference reusable, but does not by itself verify its dynamical
assumptions for another family.
